\section*{Sets and Indices}

\begin{itemize}
    \item $I$ : set of coal yards, indexed by $i$. In this instance,
    \[
    I=\{A,B\}.
    \]
    \item $J$ : set of residential areas, indexed by $j$. In this instance,
    \[
    J=\{1,2,3\}.
    \]
\end{itemize}

\section*{Parameters}

\begin{itemize}
    \item $d_{ij}$ : transportation distance from coal yard $i$ to residential area $j$ (code: \texttt{distances[(yard, area)]}).
    
    From the CSV data:
    \[
    d_{A,1}=10,\quad d_{A,2}=5,\quad d_{A,3}=6,
    \]
    \[
    d_{B,1}=4,\quad d_{B,2}=8,\quad d_{B,3}=15.
    \]
    
    \item $m_i$ : minimum monthly coal supply required from coal yard $i$ (code: \texttt{min\_supply[yard]}).
    
    For this instance:
    \[
    m_A=80,\qquad m_B=100.
    \]
    
    \item $r_j$ : monthly coal demand of residential area $j$ (code: \texttt{demand[area]}).
    
    For this instance:
    \[
    r_1=55,\qquad r_2=75,\qquad r_3=50.
    \]
\end{itemize}

\section*{Decision Variables}

\begin{itemize}
    \item $x_{ij} \ge 0$ : amount of coal shipped from coal yard $i$ to residential area $j$ in tons (code: \texttt{x[(yard, area)]}).
\end{itemize}

\section*{Objective}

Minimize total transportation cost, modeled as distance times shipped quantity:
\[
\min \sum_{i\in I}\sum_{j\in J} d_{ij}\,x_{ij}.
\]

For this instance, the objective is
\[
\min\;
10x_{A,1}+5x_{A,2}+6x_{A,3}
+4x_{B,1}+8x_{B,2}+15x_{B,3}.
\]

\section*{Constraints}

\begin{enumerate}
    \item \textbf{Minimum supply from each coal yard}
    \[
    \sum_{j\in J} x_{ij} \ge m_i
    \qquad \forall i\in I.
    \]
    Explicitly,
    \[
    x_{A,1}+x_{A,2}+x_{A,3} \ge 80,
    \]
    \[
    x_{B,1}+x_{B,2}+x_{B,3} \ge 100.
    \]

    \item \textbf{Demand satisfaction at each residential area}
    \[
    \sum_{i\in I} x_{ij} = r_j
    \qquad \forall j\in J.
    \]
    Explicitly,
    \[
    x_{A,1}+x_{B,1}=55,
    \]
    \[
    x_{A,2}+x_{B,2}=75,
    \]
    \[
    x_{A,3}+x_{B,3}=50.
    \]

    \item \textbf{Nonnegativity}
    \[
    x_{ij}\ge 0
    \qquad \forall i\in I,\; j\in J.
    \]
\end{enumerate}

\section*{Notes on Implementation}

The implemented model in \texttt{current\_heuristic.py} is a continuous linear program created with PuLP-compatible syntax and solved with \texttt{GUROBI\_CMD}. There are no integer or binary variables.

The code uses the term ``minimum supply'' for each coal yard and enforces it as a lower bound on total outbound shipments:
\[
\sum_{j} x_{ij} \ge m_i.
\]
This is the exact implemented logic, even though in a standard transportation model one might more commonly see upper-bound supply capacities instead.

Because total minimum required supply equals total demand in this instance,
\[
m_A+m_B = 80+100 = 180,
\qquad
r_1+r_2+r_3 = 55+75+50 = 180,
\]
the two coal-yard minimum-supply constraints are tight at any feasible solution, so each yard's total shipment is forced to equal its stated minimum in the solved instance.
